\begin{lemma}[5.6.11.1]
\EGAReaderTarget{ega4.c04.lem.5-6-11-1}{IV.2.s5.6.11.1-fr}
Sean $R$ un anillo, $S$ un subanillo y
$\mathfrak K=\operatorname{Ann}_S(R/S)$ el conductor de $S$ en $R$
(el mayor ideal de $S$ que es también un ideal de $R$).
\begin{enumerate}
\item[\textup{(i)}] Para todo ideal $\mathfrak J\subset\mathfrak K$ de $R$, existe una
aplicación estrictamente creciente del conjunto de los ideales $\mathfrak L$
de $S$ tales que $R.\mathfrak L=\mathfrak J$ al conjunto de los
sub-$(S/\mathfrak K)$-módulos de
$\mathfrak J/\mathfrak K\mathfrak J$.
\item[\textup{(ii)}] Si $S/\mathfrak K$ y $R$ son noetherianos y si $R$ es un
$S$-módulo de tipo finito, entonces toda sucesión creciente de ideales de $S$
contenidos en $K$ es estacionaria.
\end{enumerate}
\end{lemma}

(i) Se tiene, en efecto,
$\mathfrak K\mathfrak J=\mathfrak K(R.\mathfrak L)
=\mathfrak K\mathfrak L\subset\mathfrak L$
(pues $\mathfrak L$ es un ideal de $S$), de donde resulta la conclusión.

(ii) Sea $(\mathfrak L_n)$ una sucesión creciente de ideales de $S$ contenidos
en $\mathfrak K$. La sucesión de ideales $R.\mathfrak L_n$ de $R$ es
estacionaria puesto que $R$ es noetheriano; se puede, por tanto, suponer que
todos los ideales $R.\mathfrak L_n$ son iguales a un mismo ideal $\mathfrak J$
de $R$. Como $R$ es noetheriano, $\mathfrak J/\mathfrak K\mathfrak J$ es un
$R$-módulo de tipo finito, por tanto también un $S$-módulo de tipo finito y,
por consiguiente, un $(S/\mathfrak K)$-módulo de tipo finito. Pero como
$S/\mathfrak K$ es noetheriano, $\mathfrak J/\mathfrak K\mathfrak J$ es un
$(S/\mathfrak K)$-módulo noetheriano, y la conclusión se sigue de (i).

Vamos a aplicar este lema tomando $R=E$, $S=C$; es claro que $\mathfrak n$ es
el conductor de $C$ en $E$; además, para todo ideal $\mathfrak a$ de $C$,
$\mathfrak a/(\mathfrak n\cap\mathfrak a)$ es isomorfo a
$(\mathfrak a+\mathfrak n)/\mathfrak n$, por tanto un sub-$C$-módulo de
$C/\mathfrak n$, que es un $C$-módulo simple; basta, pues, demostrar que todo
ideal $\mathfrak a\cap\mathfrak n$ de $C$ es de tipo finito, y esto se sigue
de (5.6.11.1).

De (\textbf 0, 16.1.5) se sigue que se tiene $\dim(C)=\dim(E)=2$. Tomemos
$A=C_{\mathfrak n}$; si $U=C\setmin\mathfrak n$, el anillo de fracciones
$B=U^{-1}E$ es, por tanto, la clausura integral de $A$, y es un $A$-módulo de
tipo finito; además, como $\mathfrak m$ y $\mathfrak m'$ son los únicos
ideales primos de $E$ que contienen a $\mathfrak n$, $B$ es un anillo
semi-local cuyos anillos locales en los dos ideales maximales son isomorfos a
$E_{\mathfrak m}$ y $E_{\mathfrak m'}$, respectivamente, y son, por tanto,
de dimensiones 2 y 1, respectivamente. Esto muestra que $\dim(B)=2$, y por
tanto también $\dim(A)=2$ (\textbf 0, 16.1.5). Como $A$ es un anillo local íntegro,
es necesariamente catenario (pues todo ideal primo distinto de 0 y del ideal
maximal es necesariamente de altura 1); pero no verifica la conclusión de
(5.6.10), y \emph{a fortiori} no es universalmente catenario.

Observemos además que $\mathfrak n$ es un ideal de altura 2 en $C$, y que,
para todo ideal primo $\mathfrak p$ de altura 1 en $C$, existe \emph{un único} ideal
primo $\mathfrak p'$ de $E$ situado sobre $\mathfrak p$, necesariamente de
altura 1, y tal que $C_{\mathfrak p}=E_{\mathfrak p'}$. En efecto, existe al
menos un ideal primo $\mathfrak p'$ de $E$ situado sobre $\mathfrak p$, y se
sigue del teorema de Cohen-Seidenberg que
\oldpage[IV-2]{103}
tal ideal es necesariamente de altura 1; además, como $E$ es una
$C$-álgebra finita y $\mathfrak p\not\supset\mathfrak n$,
$C_{\mathfrak p}$ es integralmente cerrado (Bourbaki, \emph{Alg. comm.},
chap. V, § 1, n° 5, cor. 5 de la prop. 16), por tanto
$E_{\mathfrak p}=C_{\mathfrak p}$, lo que demuestra nuestras afirmaciones.

Sería interesante saber si todo anillo local íntegro noetheriano que verifica
la conclusión de (5.6.10) es universalmente catenario; esto es lo que afirmó
Nagata [33], pero su demostración no parece completa.
\end{example}
